\section{Aims, Research Questions, and Claim Boundary}
\label{ch:aims}

\subsection{Adopted end goal}

The end goal of this dissertation is to determine whether a
literature-motivated change in network dynamics can reproduce a selective
pattern of acute psilocybin-related working-memory effects more convincingly
than generic neural disruption, and to explain how that pattern emerges from
the RNN's internal dynamics. The project is not intended as a literal
pharmacological simulation of psilocybin. A task-trained RNN is used as a
controlled computational testbed for comparing candidate mechanisms.

\subsection{Required contributions}

\begin{enumerate}[leftmargin=*]
  \item Establish a credible working-memory baseline across independently
        trained checkpoints.
  \item Define a nonspecific perturbation control (independent Gaussian noise),
        retaining structured-noise work as background robustness evidence
        rather than as a direct psilocybin model.
  \item Compare literature-grounded candidate perturbations against selective
        behavioural signatures: distractor vulnerability, high-load updating
        cost, response slowing or longer settling, relative preservation under
        weaker perturbation, and stronger impairment at greater strengths.
  \item Provide a dynamical explanation in terms of delay-period decoding,
        distractor displacement and recovery, representational drift,
        memory-subspace geometry, and related hidden-state diagnostics.
\end{enumerate}

\subsection{Working research questions}

\begin{enumerate}[leftmargin=*]
  \item Which single-operator RNN perturbations, if any, reproduce the selected
        cross-task behavioural pattern against the native baseline?
  \item Does the leading candidate exceed matched-cost Gaussian disruption?
  \item What changes in recurrent population dynamics explain any selective
        behavioural match?
\end{enumerate}

\subsection{Provisional claim map}

Table~\ref{tab:claim-map} records the present evidential status. Claims that
depend on unfinished experiments remain open.

\begin{table*}[t]
\caption{Provisional claim map. Status distinguishes what the present
candidate screen supports from what still requires Sections~\ref{ch:results-specificity}
and~\ref{ch:results-dynamics}.}
\label{tab:claim-map}
\begin{tabular}{p{0.42\textwidth}p{0.22\textwidth}p{0.28\textwidth}}
\toprule
Claim & Present status & Main support / blocker \\
\midrule
Competent circular and N-back baselines exist
& Supported
& Section~\ref{ch:results-baseline}; retained checkpoints \\
Persistence $0.95$ meets the descriptive majority pattern vs native baseline
& Supported, provisional
& Section~\ref{ch:results-screen}; weak circular intervals \\
N-back load selectivity under persistence $0.95$ is robust
& Supported
& 10/10 seeds; interval excludes zero \\
Circular settling/delay/distractor effects are robust
& Not supported yet
& Seed-level intervals include zero \\
Circular and N-back hidden-state geometry is interpretable at baseline
& Supported, descriptive
& Section~\ref{ch:results-baseline}; representative seeds \\
Persistence exceeds matched-cost Gaussian disruption
& Open
& Section~\ref{ch:results-specificity} not run \\
Selective behaviour is explained by recurrent dynamics
& Open
& Section~\ref{ch:results-dynamics} not run \\
Biological identification of psilocybin mechanism
& Out of scope
& Claim boundary below \\
\bottomrule
\end{tabular}
\end{table*}

\subsection{Inferential boundaries}

The dissertation should not claim that:
\begin{itemize}[leftmargin=*]
  \item the RNN is a biological simulation of psilocybin;
  \item one perturbation strength corresponds to a human dose;
  \item behavioural similarity identifies the true neural mechanism;
  \item increased entropy alone makes a model psychedelic-like;
  \item a null specificity result is experimental failure rather than a
        scientifically informative constraint.
\end{itemize}
